\documentclass[aspectratio=169]{beamer}
\usepackage[utf8]{inputenc}
\usepackage{hyperref}
\usepackage{amsmath,amsfonts,amsthm,bm}
\usepackage{color}
\usepackage{graphicx} % Allows including images
\usepackage{booktabs} % Allows the use of \toprule, \midrule and \bottomrule in tables
\usepackage{tikz}
\usepackage{pgfplots}
\usepackage{listings}
\usepackage{courier}

\lstset{ %
  basicstyle=\scriptsize\ttfamily, % fonts that are used for the code
  breakatwhitespace=false,         % sets if automatic breaks should only happen at whitespace
  %breaklines=true,                 % sets automatic line breaking
  %captionpos=b,                    % sets the caption-position to bottom
  commentstyle=\color{gray}\textit,    % comment style
  keepspaces=true,                 % keeps spaces in text, useful for keeping indentation of code (possibly needs columns=flexible)
  keywordstyle=\color{blue},       % keyword style
  language=Python,                 % the language of the code
  %otherkeywords={*,...},          % if you want to add more keywords to the set
  rulecolor=\color{black},         % if not set, the frame-color may be changed on line-breaks within not-black text (e.g. comments (green here))
  showspaces=false,                % show spaces everywhere adding particular underscores; it overrides 'showstringspaces'
  showstringspaces=false,          % underline spaces within strings only
  showtabs=false,                  % show tabs within strings adding particular underscores
  stringstyle=\color{red}, % string literal style
  tabsize=4,	                   % sets default tabsize to 2 spaces
  columns=fixed                    % Using fixed column width (for e.g. nice alignment)
}

\hypersetup{
    colorlinks=true,
    linkcolor=red,
    filecolor=magenta,      
    urlcolor=red,
}

\DeclareMathOperator*{\argmax}{argmax}
\DeclareMathOperator*{\argmin}{argmin}
\let \vec \mathbf

\newcommand{\classname}{NANOx81}
\newcommand{\classyear}{Fall 2022}
\mode<presentation> {
    \usetheme{CambridgeUS}
    \setbeamertemplate{footline}[text line]{%
      \parbox{\linewidth}{\vspace*{-8pt}\classname\hfill\classyear\hfill\insertpagenumber}}

    %\setbeamertemplate{footline}[page number]
    \setbeamertemplate{navigation symbols}{}
}


\title[\classname Course Admin]{\classname - Data Science in Materials Science\\Course Admin}

\author{Shyue Ping Ong}
\institute[UCSD]{University of California, San Diego\\
\medskip
}
\date{\classyear} % Date, can be changed to a custom date

\begin{document}


\begin{frame}
    \titlepage % Print the title page as the first slide
\end{frame}


\begin{frame}{Course Objectives}
\huge{To provide a comprehensive introduction into the application of data science to materials science.}
\end{frame}


\begin{frame}{What will you learn in this course?}
    \begin{itemize}
        \item Enough data science (including the mathematics) to understand how to apply them to solve materials science problems.
        \item Best practices in using various data science techniques.
        \item Practical use of open-source Python packages to do data science in materials science.
    \end{itemize}
\end{frame}


\begin{frame}{What this course is \textbf{not}}
    \begin{itemize}
        \item A probability and statistics course.
        \item A replacement for a rigorous data science course.
        \item A replacement for a materials science course.
    \end{itemize}
\end{frame}


\begin{frame}{Course Plan}
    \begin{itemize}
        \item Weeks 1 and 2: Introduction to Data Science, Python and Data Wrangling
        \item Week 3: Lab 1 (Introduction to Python for Data Science and Data Wrangling)
        \item Weeks 4 and 5: Linear Methods and Unsupervised Learning
        \item Week 6: Lab 2 (Exercises in linear methods and clustering for materials science)
        \item Weeks 7 and 8: Kernel Methods, Trees and Neural Networks
        \item Weeks 9 and 10: Final Lab (Kaggle competition)
    \end{itemize}
\end{frame}


\begin{frame}{Instructors}
    \begin{itemize}
        \item Lecturer: Shyue Ping Ong (ongsp@eng.ucsd.edu)
        \item Teaching Assistant: Ji Qi (j1qi@eng.ucsd.edu)
    \end{itemize}
\end{frame}


\begin{frame}
\frametitle{Recommended Textbooks (All Free)}

\begin{columns}
\column{0.5\textwidth}
\begin{figure}
    \centering
    \includegraphics[width=0.4\textwidth]{figures/book-cover-elements-of-statistical-learning.jpg}
    \caption{The Elements of Statistical Learning: Data Mining, Inference, and Prediction, Second Edition  \href{http://www.amazon.com/dp/0387848576/ref=cm_sw_em_r_mt_dp_U_LTvbEbPHKK8VF}{[Amazon]}\href{http://web.stanford.edu/~hastie/ElemStatLearn/}{[Free PDF]}}
\end{figure}
\column{0.5\textwidth}
\begin{figure}
    \centering
    \includegraphics[width=0.4\textwidth]{figures/book-cover-python-data-science.jpg}
    \caption{Python Data Science Handbook
        \href{https://www.amazon.com/dp/1491912057/ref=cm_sw_em_r_mt_dp_U_g0vbEb9N7HVQD}{[Amazon]} \href{http://jakevdp.github.io/PythonDataScienceHandbook/}{[Free web version]}}
\end{figure}
\end{columns}
\end{frame}


\begin{frame}{Course Structure}
    \begin{itemize}
        \item Lectures/Labs (Tues/Thurs @ 1100-1220) in person.
        \item Note: Please ignore scheduled lab sessions - all labs are held in lecture times, not in a separate time.
        \item Recordings will be available online after the class.
        \item \textbf{Please bring your laptops.}
        \item Grading:
        \begin{itemize}
            \item Lab 1: 25\% - Same for NANO181 and NANO281
            \item Lab 2: 25\% - Same for NANO181 and NANO281
            \item Lab 3: 50\% - Different for NANO181 and NANO281
        \end{itemize}
    \end{itemize}
\end{frame}


\begin{frame}{Lab Assessment Criteria}
\begin{table}[]
    \centering
    \begin{tabular}{|c|c|}
    \hline
Model performance & 30\%\\
\hline
Materials Science Insights & 30\%\\
\hline
Data Science Technique & 30\%\\
\hline
Programming Style & 10\%\\
\hline
    \end{tabular}
\end{table}
    \begin{itemize}
        \item \textbf{Collaboration policy}:  Working together is highly encouraged, but each student must submit his / her own work. 
        \item To make the best use of this course, you should make sure you can do the exercises and not over-rely on your course mates.
    \end{itemize}
\end{frame}


\begin{frame}{Class etiquette}
    \begin{itemize}
        \item Interaction preferred - interruptions with questions highly encouraged.
        \item Please be punctual. Lectures will start on time.
        \item Use of laptop to follow class examples is encouraged, but please be respectful of your lecturer and classmates by not using devices for non-class applications.  All devices must be on silent mode.
    \end{itemize}
\end{frame}


\begin{frame}{Prerequisites}
    \begin{itemize}
        \item Knowledge of basic statistics (e.g., Gaussian distributions, Bayes theorem, etc.)
        \item Knowledge of basic linear algebra (e.g., matrix multiplications, eigenvalue decompositions, inverses)
        \item Some programming experience. Ideally, experience in the Python programming language would be helpful.
        \item 1st Homework (ungraded):
        \begin{enumerate}
            \item Go to \href{https://colab.research.google.com/}{Google Colab}.
            \item Create a new notebook.
            \item Go through items 1-3 in the \href{http://docs.python.org/3/tutorial/}{official Python tutorial} – please run through the actual tutorial line by line. It should not take you more than 30 mins to do the whole thing.
            \item Extra: Briefly read through item 4 in the tutorial on flow control (if and for statements, especially).
        \end{enumerate}
    \end{itemize}
\end{frame}

\begin{frame}{Course Admin}
    \begin{itemize}
        \item Canvas for all course admin, including announcements/communications, lecture slides, submission of labs, useful resources (e.g., ebooks, websites, etc.)
        \item \href{https://github.com/materialsvirtuallab/nano281}{NANOx81 Github Page} - all labs with instructions.
        \begin{figure}
            \centering
            \includegraphics[width=0.25\textwidth]{figures/QR-NANO281Github.png}
            \label{fig:my_label}
        \end{figure}
    \end{itemize}
\end{frame}


\begin{frame}{Questions and Feedback}
    \begin{itemize}
        \item Questions welcomed at any time during or after lectures
        \item \classname is very new - the instructors will try their best, but I would ask for you to be tolerant of any issues while we continue to improve the curriculum and labs.
        \item Your \textbf{feedback} is invaluable in shaping the current course as well as future courses. 
        \item Email all feedback directly to \href{mailto:ongsp@eng.ucsd.edu}{ongsp@ucsd.edu}.
    \end{itemize}
\end{frame}


\begin{frame}
    \Huge{\centerline{The End}}
\end{frame}

\end{document}

